
This paper investigates formal repair methods for LLM-generated Python code by integrating three approaches -- grammar-constrained decoding (SynCode), SMT solving (Z3), and iterative type checking (mypy) -- into a unified Python-native pipeline evaluated with CodeLlama-7B on HumanEval. A Failure Mode Distribution (FMD) analysis is applied to characterize which repair channels can activate for this model-benchmark pair. The principal finding is a null result: across 2,680 completions from 134 problems, CodeLlama-7B produces zero type errors, rendering the mypy repair channel inapplicable. This outcome indicates that the precondition assumed by the iterative type-checking repair literature -- the presence of type errors in generated code -- is not met for this model-benchmark combination. Additionally, 97.5\% of observed failures are syntax-level errors. SynCode reduces AST parse failures by 7.5 percentage points (delta_ast = 0.075), though this reduction does not reach statistical significance at N = 20 problems (bootstrap 95\% CI: [-0.025, 0.220], p = 0.1186). Z3 constraint encoding is applicable to 33\% of HumanEval problems (54/164). The planned SynCode-Z3 complementarity experiment was not executed in this pipeline run; its result remains inconclusive.
